\section*{Bab \ref{ch:g3:mult} --- Perkalian}

\begin{solution}{exo:g3:mult:1}
$8 + 8 + 8 + 8 = 4 \times 8 = 32$; \quad
$5 + 5 + 5 = 3 \times 5 = 15$; \quad
$20 \times 5 = 100$.
\end{solution}

\begin{solution}{exo:g3:mult:2}
$42$; \quad $32$; \quad $81$; \quad $56$; \quad $36$.
\end{solution}

\begin{solution}{exo:g3:mult:3}
Kedua persegi panjang memuat $24$ titik yang sama --- yang satu
adalah yang lain yang dimiringkan: $4 \times 6 = 6 \times 4$.
\end{solution}

\begin{solution}{exo:g3:mult:4}
$560$; \quad $800$; \quad $0$; \quad $45$; \quad $300$.
\end{solution}

\begin{solution}{exo:g3:mult:5}
$132 \times 3 = 396$; \quad
$217 \times 4 = 868$; \quad
$408 \times 7 = 2\,856$.
\end{solution}

\begin{solution}{exo:g3:mult:6}
$31 \times 5 = 30 \times 5 + 5 = 155$; \quad
$42 \times 3 = 120 + 6 = 126$; \quad
$25 \times 6 = 20 \times 6 + 5 \times 6 = 120 + 30 = 150$.
\end{solution}

\begin{solution}{exo:g3:mult:7}
$99 \times 7 = 700 - 7 = 693$; \quad
$101 \times 6 = 600 + 6 = 606$; \quad
$98 \times 5 = 500 - 10 = 490$.
\end{solution}

\begin{solution}{exo:g3:mult:8}
$389 \times 5$: taksiran $400 \times 5 = 2\,000$; tepatnya
$1\,945$.

$612 \times 8$: taksiran $600 \times 8 = 4\,800$; tepatnya $4\,896$.
\end{solution}

\begin{solution}{exo:g3:mult:9}
$38 \times 6 = 228$. Ada $228$ telur di dalam kotak itu.
\end{solution}

\begin{solution}{exo:g3:mult:10}
Meja: $8 \times 4 = 32$. Murid: $32 \times 2 = 64$. Ruangan itu
dapat menampung $64$ murid.
\end{solution}

\begin{solution}{exo:g3:mult:11}
$24$ muncul dalam persegi tabel sebagai $1 \times 24$, $2 \times 12$,
$3 \times 8$, $4 \times 6$ --- dan pasangan yang sama dibalik.
Susunan yang mungkin: $1$ baris berisi $24$, $2$ baris berisi $12$,
$3$ baris berisi $8$, $4$ baris berisi $6$, $6$ baris berisi $4$,
$8$ baris berisi $3$, $12$ baris berisi $2$, $24$ baris berisi $1$.
\end{solution}

\begin{solution}{exo:g3:mult:12}
Zahra memecah $47$ menjadi $40 + 7$ (cara bersusun di dalam
kepalanya); Lia mengganti $47$ dengan $50$ yang ramah lalu
membetulkannya dengan mengambil $3 \times 6$ yang terhitung
berlebih. Untuk $38 \times 4$: dengan cara Zahra,
$120 + 32 = 152$; dengan cara Lia, $40 \times 4 - 2 \times 4 = 160 - 8 =
152$. Keduanya berhasil --- cara Lia lebih cepat di sini karena $38$
dekat dengan $40$.

\end{solution}
